\section{From Action Commands to Motion and Task Outcomes}
\label{sec:motion}
This section follows the last two links. Do changed commands move the robot, and does the size of that motion tell how much success is lost?

\begin{table}[t]
\centering\small
\setlength{\tabcolsep}{2.8pt}
\begin{tabular}{@{}lrrrrrrr@{}}
\toprule
 & Motion & \multicolumn{2}{c}{Success} & \multicolumn{3}{c}{Perturbation} & Rerun\\
\cmidrule(lr){3-4}\cmidrule(lr){5-7}
Family & ratio & clean & pert. & lost & gained & all & all\\
\midrule
\multicolumn{8}{@{}l}{$\pi_{0.5}$}\\
Background & 0.93 & 94.6 & 94.6 & 3.6 & 3.6 & 7.1 & 3.6\\
Camera & 2.68 & 94.6 & 74.1 & 25.0 & 4.5 & 29.5 & 3.6\\
Light$^\dagger$ & 0.57 & 91.1 & -- & -- & -- & -- & --\\
Sensor Noise & 1.32 & 92.9 & 84.8 & 14.3 & 6.3 & 20.5 & 7.1\\
\multicolumn{8}{@{}l}{$\pi_0$}\\
Background & 0.62 & 89.3 & 85.7 & 7.1 & 3.6 & 10.7 & 7.1\\
Camera & 1.03 & 86.6 & 66.1 & 22.3 & 1.8 & 24.1 & 5.4\\
Light$^\dagger$ & 0.24 & 87.5 & -- & -- & -- & -- & --\\
Sensor Noise & 0.58 & 92.9 & 78.6 & 16.1 & 1.8 & 17.9 & 0.0\\
\bottomrule
\end{tabular}
\caption{\textbf{How far each perturbation moves the robot, and what it does to task outcomes.} Motion ratio: Sec.~\ref{sec:sensitivity}. The other columns are percentages over the 28 common tasks: success of the clean and the perturbed run; outcomes that succeed when clean and fail when perturbed (lost), the reverse (gained), and both (all); and outcomes that differ between the two clean runs (rerun). $^\dagger$Some Light outcomes were not recorded, so no rates are given; Fig.~\ref{fig:behavior} shows the recorded ones (bounds: Supplement, Sec.~S6).}
\label{tab:outcomes}
\end{table}


\subsection{Command Changes Reach the Robot}
\label{sec:early}
Section~\ref{sec:restoration} judged restoration by its effect on the first five commands, so we first check that such changes reach the robot. For every pair, on all 40 tasks, we compare how much the first five commands changed with how far apart the two end effectors are after 5 controls and after 60.

Early on, the two agree almost perfectly: the rank correlation is 0.95 in $\pi_{0.5}$ and 0.97 in $\pi_0$. This is expected, since the first five controls simply execute those commands. After 60 controls, eleven replans later, the correlation is 0.47 and 0.41: weaker, but still clearly present (Supplement, Sec.~S5.1). A perturbation's first effect on the commands thus remains visible in the robot's motion well into the episode.

The command change is also systematic: with a different draw of policy noise, the same perturbation changes the commands in nearly the same way (agreement 0.88--0.92 in $\pi_{0.5}$ and 0.80--0.89 in $\pi_0$, where 1 is identical; Supplement, Sec.~S7). The first five commands are therefore a repeatable readout of what a perturbation does to the robot.

\subsection{The Size of a Motion Change Is Not Its Cost}
\label{sec:sensitivity}
How large is a motion change? Two clean runs already differ, because each draws its own policy noise. We use that difference as the yardstick (Sec.~\ref{sec:comparisons}). The \emph{motion ratio} is the distance between a perturbed run and its clean run, divided by the distance between two clean runs, after five controls. Below 1, the perturbation moves the robot less than a clean rerun does. Table~\ref{tab:outcomes} lists it with the outcomes.

\textbf{A small motion change can cost success.} Take $\pi_0$ under Sensor Noise. After five controls, the perturbed arm is 3.5\,mm from where the clean run put it, while a clean rerun lands 6.1\,mm away. Yet in these settings a rerun never changes whether the task succeeds, whereas the perturbation costs 14.3 points of success (interval $[-28.6,-2.7]$; Fig.~\ref{fig:behavior}a).

\textbf{Motion size is a poor predictor of lost success.} Across families, larger motion changes mostly come with larger losses: Camera moves the robot most and costs the most success in both policies (Table~\ref{tab:outcomes}). For individual runs the link is weak. We rank the pairs by their motion change and ask how well the ranking separates pairs that lose a success from pairs that keep it. This is the AUROC (0.5 is chance, 1 is perfect), computed post hoc. With all families pooled it is 0.72 in $\pi_{0.5}$ and 0.64 in $\pi_0$; after 60 controls it is 0.60 and 0.68; and raw distance or the size of the command change gives nearly the same values as the ratio. The pooled values mostly reflect differences between families: within a family the AUROC drops to 0.62 and 0.55, and for $\pi_0$ under Sensor Noise to 0.49 (Fig.~\ref{fig:behavior}c; Supplement, Secs.~S5 and~S6).

\textbf{Three further comparisons show the same gap.} Under Camera, $\pi_0$ moves about as far as a clean rerun (1.03$\times$) and loses 20.5 points (Table~\ref{tab:outcomes}). In a study of ten visual conditions for $\pi_{0.5}$ on original LIBERO, conditions that change the commands more lose more success on average (rank correlation $-0.90$ across the ten conditions), but task by task the link all but disappears ($-0.075$ across the 400 combinations of task and condition; Supplement, Sec.~S9.1): the same pattern as between and within families. And removing a headlight changes the first five commands 5.1 times as much as a dimming matched in image statistics, yet both succeed in 77 of 80 cases (Supplement, Sec.~S9.2).

\textbf{Net success hides lost and gained successes.} Under a perturbation some clean successes fail and some clean failures succeed; the net change shows only their difference. Under Sensor Noise, $\pi_{0.5}$ gains nearly half as many successes as it loses, so its net loss of 8.0 points understates the 20.5\% of outcomes that change (Table~\ref{tab:outcomes}; Fig.~\ref{fig:behavior}b). A clean rerun alone changes up to 7.1\% of outcomes (10.7\% for $\pi_{0.5}$ under Light), so changes must be read against it: lost successes exceed the rerun level under Camera and Sensor Noise in both policies, while gained successes and the changes under Background do not (Supplement, Sec.~S6.1). Paired runs under combined perturbations show such cancellations~\cite{sawada2026paired}; the rerun adds the baseline for reading them.

The gap between sensitivity and fragility thus opens at the last link: a perturbation's effect can be followed from shared feature directions to commands and into the robot's motion, but its size says little about which runs fail. Robustness has to be judged by paired outcomes against clean reruns.
